\documentclass[11pt]{article}

\usepackage{xeCJK}
\setCJKmainfont{STSong}
\setCJKsansfont{Heiti SC}

\input{../../templates/template.LaTeX.homework/structure.tex}

%----------------------------------------------------------------------------------------
%  ASSIGNMENT INFORMATION
%----------------------------------------------------------------------------------------
\newcommand{\assignmentQuestionName}{习题}
\newcommand{\assignmentClass}{离散数学}
\newcommand{\assignmentTitle}{第4章 基数与可数集 — 第12次作业}
\newcommand{\assignmentAuthorName}{乐绎华}
\newcommand{\studentID}{23363017}
\newcommand{\assignmentClassInstructor}{左孝凌、刘永才《离散数学》}
\newcommand{\assignmentDueDate}{}

\begin{document}

\maketitle
\thispagestyle{empty}
\newpage

%----------------------------------------------------------------------------------------
%  4-4(1d) 构造双射函数
%----------------------------------------------------------------------------------------

\begin{question}{4-4(1d) 构造双射函数}
  构造一个从 $A = R$ 到 $B = (0, \infty)$ 的双射函数，说明 $A$ 和 $B$ 具有相同的势。

  \begin{answer}
    \textbf{构造函数：}
    定义函数 $f: R \to (0, \infty)$ 如下：
    $$f(x) = e^x$$

    \textbf{证明 $f$ 是双射：}
    \begin{enumerate}
      \item \textbf{单射性：} 对于任意 $x_1, x_2 \in R$，若 $f(x_1) = f(x_2)$，即 $e^{x_1} = e^{x_2}$。两边取自然对数得 $\ln(e^{x_1}) = \ln(e^{x_2})$，即 $x_1 = x_2$。故 $f$ 是单射。
      \item \textbf{满射性：} 对于任意 $y \in (0, \infty)$，令 $x = \ln y$。因为 $y > 0$，则 $x \in R$ 且 $f(x) = e^{\ln y} = y$。故 $f$ 是满射。
    \end{enumerate}

    综上所述，$f$ 是从 $R$ 到 $(0, \infty)$ 的双射函数。

    \begin{info}[结论:]
      由于存在双射函数，$A = R$ 与 $B = (0, \infty)$ 等势，即 $K[R] = K[(0, \infty)] = \aleph$。
    \end{info}
  \end{answer}
\end{question}

%----------------------------------------------------------------------------------------
%  4-5(1b, 1c) 集合的势
%----------------------------------------------------------------------------------------

\begin{question}{4-5(1b, 1c) 集合的势}
  下列集合 $A$ 的势是什么？
  \begin{enumerate}
    \item[b)] $A = \{\langle p, q \rangle \mid p, q \text{ 都是有理数}\}$
    \item[c)] $A$ 是由所有半径为1，圆心在 $x$ 轴上的圆周所组成的集合
  \end{enumerate}

  \begin{answer}
    \textbf{b) $A = \{\langle p, q \rangle \mid p, q \in Q\}$ 的势}

    由题意知 $A = Q \times Q$。
    \begin{itemize}
      \item 已知有理数集 $Q$ 是可数集，即 $K[Q] = \aleph_0$。
      \item 根据定理：两个可数集的笛卡尔积仍然是可数集。
      \item 因此 $Q \times Q$ 是可数集。
    \end{itemize}
    故 $K[A] = \aleph_0$。

    \vspace{1em}
    \textbf{c) 半径为1，圆心在 $x$ 轴上的圆周集合的势}

    每一个此类圆周由其圆心唯一确定。圆心在 $x$ 轴上，其坐标形式为 $(x, 0)$，其中 $x \in R$。
    设 $S$ 为此类圆周的集合，构造映射 $f: S \to R$：
    $$f(\text{圆心为 } (x, 0) \text{ 的圆周}) = x$$
    显然 $f$ 是双射。
    因为 $K[R] = \aleph$，故 $K[A] = \aleph$。

    \begin{info}[总结:]
      b) 的势为 $\aleph_0$（可数集）；c) 的势为 $\aleph$（连续统势）。
    \end{info}
  \end{answer}
\end{question}

%----------------------------------------------------------------------------------------
%  4-5(5) 笛卡尔积的可数性
%----------------------------------------------------------------------------------------

\begin{question}{4-5(5) 笛卡尔积的可数性}
  证明：有限集 $A$ 和可数集 $B$ 的笛卡尔积集 $A \times B$ 是可数集。

  \begin{answer}
    \begin{proof}
      若 $A$ 为空集，则 $A \times B = \emptyset$，是有穷集（至多可数集）。

      若 $A$ 为非空有限集，设 $A = \{a_1, a_2, \dots, a_n\}$，其中 $n = |A| \in Z^+$。
      则 $A \times B$ 可以表示为：
      $$A \times B = \bigcup_{i=1}^n (\{a_i\} \times B)$$

      对于每一个集合 $B_i = \{a_i\} \times B$：
      定义映射 $g_i: B \to B_i$ 为 $g_i(b) = \langle a_i, b \rangle$。
      容易验证 $g_i$ 是双射，故 $B_i \sim B$。
      由于 $B$ 是可数集，故每个 $B_i$ 也是可数集。

      由定理“有限个可数集的并集是可数集”可知：
      $A \times B = \bigcup_{i=1}^n B_i$ 是可数集。
    \end{proof}

    \begin{info}[注:]
      本题结论表明，可数性的性质在有限集合的笛卡尔积运算下是保持的。
    \end{info}
  \end{answer}
\end{question}

\end{document}
